\chapter*{Introduction G\'{e}n\'{e}rale}
\addcontentsline{toc}{chapter}{Introduction G\'{e}n\'{e}rale}

La transformation num\'{e}rique du commerce ext\'{e}rieur constitue un levier strat\'{e}gique majeur pour la comp\'{e}titivit\'{e} \'{e}conomique du Maroc. Depuis la cr\'{e}ation du Guichet Unique National en 2008, \textbf{PortNet S.A.} joue un r\^{o}le central d'interconnexion entre les op\'{e}rateurs \'{e}conomiques (importateurs, exportateurs, transitaires), les administrations publiques (Douanes, OMPIC, DGI, ONSSA) et les \'{e}tablissements bancaires. La plateforme rassemble aujourd'hui plus de 50\,000 usagers actifs et constitue le passage oblig\'{e} de toute op\'{e}ration de d\'{e}douanement au Maroc.

L'acc\`{e}s aux services informatiques de PortNet et \`{a} la plateforme BADR de l'Administration des Douanes et Imp\^{o}ts Indirects (ADII) est conditionn\'{e} par une proc\'{e}dure de souscription r\'{e}glementaire. Chaque op\'{e}rateur doit fournir un dossier compos\'{e} de neuf pi\`{e}ces justificatives : le Mod\`{e}le~J du Registre de Commerce, le Bulletin de Notification de l'Identifiant Fiscal, l'Attestation de Taxe Professionnelle, le Certificat de l'Identifiant Commun de l'Entreprise (ICE), l'Attestation de Domiciliation, l'Attestation de Relev\'{e} d'Identit\'{e} Bancaire (RIB), la Carte Nationale d'Identit\'{e} \'{E}lectronique (CNIE), les Statuts de la Soci\'{e}t\'{e} et, le cas \'{e}ch\'{e}ant, une Procuration.

Dans le processus actuel, l'instruction de chaque dossier repose sur un contr\^{o}le visuel manuel r\'{e}alis\'{e} par des agents v\'{e}rificateurs. Cette approche pr\'{e}sente plusieurs limitations structurelles : des d\'{e}lais d'instruction pouvant atteindre 48 \`{a} 72 heures ouvr\'{e}es en p\'{e}riode de pointe, un taux d'erreur humaine non n\'{e}gligeable face \`{a} la diversit\'{e} des formats de documents administratifs marocains, et l'impossibilit\'{e} pratique de r\'{e}aliser manuellement certaines v\'{e}rifications math\'{e}matiques critiques, telles que le calcul des cl\'{e}s de contr\^{o}le Modulo~97 pour les identifiants ICE et les num\'{e}ros de compte RIB.

Ces difficult\'{e}s sont amplifi\'{e}es par la sp\'{e}cificit\'{e} des documents administratifs marocains. Les pi\`{e}ces justificatives partagent souvent des mentions l\'{e}gales identiques (\og Royaume du Maroc \fg{}, num\'{e}ros ICE, RC, IF) qui rendent la discrimination entre types de documents particuli\`{e}rement ardue pour un syst\`{e}me automatis\'{e}. Par ailleurs, les formulaires officiels de la DGI et de l'OMPIC pr\'{e}sentent des agencements de texte non standards o\`{u} les valeurs apparaissent avant leurs libell\'{e}s, pi\'{e}geant les algorithmes d'extraction conventionnels. Enfin, les fautes de frappe et transpositions sur les identifiants des g\'{e}rants (num\'{e}ros CIN) n\'{e}cessitent un m\'{e}canisme de tol\'{e}rance qui pr\'{e}serve la rigueur du contr\^{o}le tout en \'{e}vitant les rejets injustifi\'{e}s.

C'est dans ce contexte que s'inscrit le pr\'{e}sent Projet de Fin d'Ann\'{e}e (PFA), r\'{e}alis\'{e} au sein du D\'{e}partement D\'{e}veloppement et Ing\'{e}nierie Logicielle de PortNet S.A. L'objectif est de concevoir et d'impl\'{e}menter un syst\`{e}me intelligent de v\'{e}rification documentaire capable d'automatiser l'int\'{e}gralit\'{e} du processus de contr\^{o}le : de l'extraction du texte brut des fichiers PDF \`{a} la validation crois\'{e}e des donn\'{e}es entre tous les documents du dossier, en passant par la classification par emplacement et la v\'{e}rification des sommes de contr\^{o}le math\'{e}matiques.

Le syst\`{e}me r\'{e}alis\'{e} repose sur une architecture en pipeline \`{a} trois niveaux. Le premier niveau assure l'ingestion des fichiers PDF par extraction textuelle native (PyMuPDF) avec repli OCR pour les documents num\'{e}ris\'{e}s. Le deuxi\`{e}me niveau impl\'{e}mente un classifier \`{a} deux couches qui combine des ancres textuelles positives propres \`{a} chaque type de document avec des exclusions n\'{e}gatives pour rejeter les faux positifs inter-emplacements. Le troisi\`{e}me niveau orchestre les validations : extraction des variables m\'{e}tier par des parsers regex d\'{e}di\'{e}s, contr\^{o}le des cl\'{e}s Modulo~97, v\'{e}rification de la r\'{e}cence, correspondance floue des cha\^{\i}nes de caract\`{e}res et ex\'{e}cution de la matrice de coh\'{e}rence crois\'{e}e inter-documents.

L'architecture technique est fond\'{e}e sur un backend FastAPI avec persistance PostgreSQL, un frontend mono-page (SPA) reproduisant fid\`{e}lement l'ergonomie du portail PortNet, et une couche d'extraction documentaire en Python pur. Le syst\`{e}me a \'{e}t\'{e} valid\'{e} par une campagne de tests automatiques couvrant 64 cas de test avec un taux de r\'{e}ussite de 100\%.

Le pr\'{e}sent rapport est organis\'{e} en cinq chapitres. Le Chapitre~1 introduit l'organisme d'accueil, le contexte du projet, l'\'{e}tude et la critique de l'existant, la probl\'{e}matique, le cahier des charges fonctionnel et non-fonctionnel, ainsi que la m\'{e}thodologie de conduite de projet. Le Chapitre~2 pr\'{e}sente l'\'{e}tat de l'art concernant l'extraction textuelle, la reconnaissance optique de caract\`{e}res, la classification de documents, la validation math\'{e}matique et la correspondance floue. Le Chapitre~3 d\'{e}crit la conception et la mod\'{e}lisation UML du syst\`{e}me : cas d'utilisation, s\'{e}quences, activit\'{e}s, \'{e}tats, architecture et classes. Le Chapitre~4 d\'{e}taille la r\'{e}alisation technique, le pipeline d'impl\'{e}mentation, les extracteurs, les validateurs et l'interface utilisateur. Enfin, le Chapitre~5 pr\'{e}sente les tests, les r\'{e}sultats exp\'{e}rimentaux et l'\'{e}valuation des performances du syst\`{e}me.
